\chapterimage{figs/chapterhead.png}
\chapter{Packaging and CI}
\label{chap:packaging-and-ci}

This chapter is based on the current \texttt{WorkInProgress} branch.
The recorded baseline commit is \texttt{0af9b8ec8b3961dd4cde05e690fc250cccd22ab0}
(from pull request \#522, \texttt{packaging: integrate Launcher into Debian
installation}).
It uses the active build and packaging manifests in
\texttt{CMakePresets.json}, \texttt{debian/}, \texttt{packaging/linux/}, and
\texttt{.github/workflows/}. It documents the present build, CI, and Debian
packaging surface as it exists now, not a future release plan.

\noindent\textbf{Objective.} Describe the current CI topology, Debian package
layout, and packaging boundaries that are available in the repository today.

\noindent\textbf{Audience.} Developers, maintainers, and reviewers who need to
inspect or reproduce the current build, packaging, or pipeline behavior.

\noindent\textbf{Scope.} Cover the ordinary CI workflow, the Debian package
workflow, the generated package names, the install layout, the Launcher and
Dispatcher packaging contract, the AppStream and lintian checks, the package
lifecycle validation job, the artifact outputs, and the current release
boundary.

\noindent\textbf{Status.} Confirmed by source inspection and by a local and
GitHub Actions execution of the full package build and lifecycle validation
on 2026-08-20. The chapter records the current automation surface rather than
a hypothetical future gate.

The chapter follows this sequence:
Figure~\ref{fig:packaging-ci-topology},
Figure~\ref{fig:packaging-debian-pipeline},
Figure~\ref{fig:packaging-artifact-flow},
Figure~\ref{fig:packaging-package-layout}, and
Figure~\ref{fig:packaging-release-boundary}.

\section{Current CI topology}
\label{sec:packaging-current-ci-topology}

The ordinary CI workflow is \texttt{.github/workflows/genesys-ci.yml}. It
currently:
\begin{itemize}
\item runs on \texttt{ubuntu-24.04};
\item uses \texttt{QT\_QPA\_PLATFORM=offscreen};
\item configures, builds, and tests the \texttt{tests-unit} preset;
\item installs the current build dependencies needed by the workflow;
\item runs the GUI GMDD diagnostics job as a separate check;
\item uploads the GUI GMDD diagnostics artifact on failure or success.
\end{itemize}

Its pull request trigger set currently includes:
\begin{itemize}
\item \texttt{20261};
\item \texttt{WorkInProgress};
\item \texttt{WiP20261};
\item \texttt{currentStable};
\item \texttt{master}.
\end{itemize}
It also supports \texttt{workflow\_dispatch}.

The GUI GMDD diagnostics job is intentionally narrower than the main test job.
It configures the \texttt{tests-unit} preset, builds
\texttt{genesys\_test\_gui\_gmdd\_layout}, and runs a focused gtest filter for
the layout regression cases that are currently recorded in the workflow.

\begin{figure}[htbp]
\centering
% FIGURE-SPEC-BEGIN
% ID: fig:packaging-ci-topology
% Source of truth:
%   - ../.github/workflows/genesys-ci.yml
%   - ../README.md
%   - ../docs/ai_assistants/STATUS.md
%   - ../CMakePresets.json
% Purpose:
%   Show the current CI topology, triggers, environment, and the split between
%   the main tests-unit job and the GUI GMDD diagnostics job.
% Required visual content:
%   - PR/workflow_dispatch triggers
%   - ubuntu-24.04 runner
%   - offscreen Qt environment
%   - tests-unit configure/build/test path
%   - GUI GMDD diagnostics path
% Accessibility description:
%   A pipeline diagram that makes the current CI jobs and their claims explicit.
% Validation:
%   The workflow names, runner, presets, and trigger branches must match the
%   current YAML.
% FIGURE-SPEC-END
\figureplaceholder{figs/generated/packaging-ci-topology}
\caption{Current CI topology and claim boundaries.}
\label{fig:packaging-ci-topology}
\end{figure}

Figure~\ref{fig:packaging-ci-topology} should make the claim split explicit.
The ordinary CI job validates the main regression path, while the GUI GMDD job
only confirms the narrow layout-regression surface that it targets.

\section{Debian packaging workflow}
\label{sec:packaging-debian-workflow}

The Debian package workflow is \texttt{.github/workflows/genesys-debian-package.yml}.
It runs on \texttt{ubuntu-24.04} and is split into two jobs.

\noindent\textbf{Triggers.} The workflow supports \texttt{workflow\_dispatch}
and triggers on both push and pull-request events targeting
\texttt{WorkInProgress}, \texttt{20262}, \texttt{currentStable}, and
\texttt{master}, restricted to path filters covering \texttt{debian/**},
\texttt{packaging/linux/**}, \texttt{source/**},
\texttt{scripts/run-debian-package-workflow.sh}, \texttt{CMakeLists.txt},
\texttt{CMakePresets.json}, and the workflow file itself. Because
\texttt{WorkInProgress} is now a trigger branch (unlike the historical
\texttt{WiP20261}/\texttt{20261} filter), an ordinary pull request against
\texttt{WorkInProgress} that touches packaging-relevant paths runs this
workflow automatically, and a push to \texttt{WorkInProgress} after merge
runs it again on the merge commit.

\noindent\textbf{Job 1 -- \texttt{Build Debian packages}.} It:
\begin{itemize}
\item records provenance (source head SHA, checkout SHA, toolchain versions);
\item installs the packaging toolchain: \texttt{debhelper}, \texttt{devscripts},
  \texttt{dpkg-dev}, \texttt{cmake}, \texttt{ninja-build}, \texttt{g++},
  \texttt{pkgconf}, \texttt{qt6-base-dev(-tools)}, \texttt{libgl-dev},
  \texttt{lintian}, \texttt{appstream}, \texttt{gpgv}, \texttt{zstd};
\item runs \texttt{dpkg-buildpackage -us -uc -b};
\item collects the build log, \texttt{.buildinfo}/\texttt{.changes}, and the
  produced \texttt{.deb} files;
\item verifies that exactly the five expected binary packages were produced,
  one each;
\item validates
  \texttt{packaging/linux/io.github.rlcancian.genesys.metainfo.xml} with
  \texttt{appstreamcli validate --no-net};
\item runs \texttt{lintian --fail-on error} over all five \texttt{.deb} files,
  so any Lintian \texttt{error}-severity finding fails the job while warnings
  are recorded but do not block it;
\item uploads the \texttt{genesys-debian-packages} and
  \texttt{genesys-debian-diagnostics} artifacts.
\end{itemize}

\noindent\textbf{Job 2 -- \texttt{Validate Debian package lifecycle}.} It
depends on job 1 and only runs if the build job succeeds. It downloads the
exact \texttt{genesys-debian-packages} artifact produced by job 1 (never a
different run's artifact), installs \texttt{xvfb}, \texttt{curl},
\texttt{iproute2}, and \texttt{python3}, then runs
\texttt{packaging/linux/validate-debian-lifecycle.sh}, uploading its output
as the \texttt{genesys-debian-lifecycle-evidence} artifact. Section
\ref{sec:packaging-lifecycle-validation} describes what that script exercises.

The package workflow is separate from ordinary regression CI. A green
\texttt{tests-unit} run does not imply package lifecycle validation. A green
build job does, however, unlock the lifecycle job in the same workflow run;
only both jobs succeeding together constitute full packaging evidence, and
neither replaces the ordinary regression workflow.

\begin{figure}[htbp]
\centering
% FIGURE-SPEC-BEGIN
% ID: fig:packaging-debian-pipeline
% Source of truth:
%   - ../.github/workflows/genesys-debian-package.yml
%   - ../debian/rules
%   - ../debian/control
%   - ../packaging/linux/io.github.rlcancian.genesys.metainfo.xml
%   - ../packaging/linux/validate-debian-lifecycle.sh
% Purpose:
%   Show the current two-job Debian package pipeline: build/validate-metadata,
%   then a dependent lifecycle-validation job that consumes the exact same
%   run's package artifact.
% Required visual content:
%   - workflow_dispatch and path-triggered entry (push/PR to WorkInProgress,
%     20262, currentStable, master)
%   - dpkg-buildpackage
%   - AppStream validation
%   - lintian --fail-on error
%   - package verification and upload (Job 1)
%   - lifecycle job depending on Job 1, downloading its artifact
% Accessibility description:
%   A left-to-right pipeline that shows how source checkout becomes package
%   artifacts and diagnostics, then feeds a dependent lifecycle-validation job.
% Validation:
%   The commands, job names, and filenames must match the current workflow
%   and packaging files.
% FIGURE-SPEC-END
\figureplaceholder{figs/generated/packaging-debian-pipeline}
\caption{Current Debian package build and lifecycle-validation pipeline.}
\label{fig:packaging-debian-pipeline}
\end{figure}

Figure~\ref{fig:packaging-debian-pipeline} should show the Debian packaging path
as a distinct, two-job pipeline rather than as a side effect of ordinary CI.

\section{Package layout and install surface}
\label{sec:packaging-package-layout}

The Debian source package is \texttt{genesys-simulator}. The binary packages
currently declared in \texttt{debian/control} are:
\begin{itemize}
\item \texttt{genesys-common} -- the stable launcher, the non-interactive
  dispatcher, \texttt{/etc/genesys/update.conf}, and the public
  \texttt{genesys-mcp} dispatcher entry point;
\item \texttt{genesys-shell} -- the public \texttt{genesys-shell} wrapper and
  the Debian-provided system fallback shell;
\item \texttt{genesys-worker} -- the public \texttt{genesys-worker} wrapper,
  the legacy \texttt{genesys-web} command alias, and the system fallback
  worker;
\item \texttt{genesys-gui} -- the Qt desktop GUI wrapper, the system fallback
  GUI, and the desktop integration files; it depends on \texttt{genesys-common},
  \texttt{genesys-shell}, and \texttt{genesys-worker} so that
  \texttt{apt install genesys-gui} alone yields a complete supported base
  install;
\item \texttt{genesys-web} -- a transitional package with no payload of its
  own, depending on \texttt{genesys-worker}, kept only so upgrades from the
  historical \texttt{genesys-web} package name continue to work.
\end{itemize}

Every non-\texttt{genesys-common} package depends on \texttt{genesys-common}
at the exact same binary version (\texttt{genesys-common (= \$\{binary:Version\})}),
and \texttt{genesys-worker} carries \texttt{Breaks}/\texttt{Replaces} on older
\texttt{genesys-web} versions so ownership of the \texttt{genesys-web} command
transfers cleanly.

The current \texttt{debian/rules} file builds the package with
\texttt{cmake+ninja}, uses \texttt{build-debian} as the build directory, sets
\texttt{CMAKE\_BUILD\_TYPE=Release}, installs under \texttt{/usr}, enables the
GUI, worker, terminal (shell), and Launcher application paths, and disables
the test targets and Python integration for the Debian build; it also enables
full hardening (\texttt{DEB\_BUILD\_MAINT\_OPTIONS=hardening=+all}).

The install manifests currently place:
\begin{itemize}
\item \texttt{/usr/bin/genesys-gui}, \texttt{/usr/bin/genesys-shell},
  \texttt{/usr/bin/genesys-worker}, \texttt{/usr/bin/genesys-web},
  \texttt{/usr/bin/genesys-mcp} -- stable public wrapper commands;
\item \texttt{/usr/libexec/genesys/genesys-launcher} and
  \texttt{/usr/libexec/genesys/genesys-dispatch} -- the two mechanisms those
  wrappers invoke;
\item \texttt{/usr/libexec/genesys/system/bin/\{genesys-gui,genesys-shell,genesys-worker\}}
  -- the Debian-provided system fallback executables;
\item \texttt{/etc/genesys/update.conf} -- the administrative update policy
  conffile;
\item \texttt{/usr/share/man/man1/*.1} -- one manual page per public command;
\item \texttt{/usr/share/applications/io.github.rlcancian.genesys.desktop};
\item \texttt{/usr/share/metainfo/io.github.rlcancian.genesys.metainfo.xml};
\item \texttt{/usr/share/icons/hicolor/scalable/apps/io.github.rlcancian.genesys.svg}.
\end{itemize}

The desktop entry launches \texttt{genesys-gui} and marks the application as a
non-terminal desktop app. The metainfo file describes the desktop package and
its launchable desktop ID. Section~\ref{sec:packaging-launcher-contract}
covers the wrapper/launcher/dispatcher/system-fallback contract in detail.

\begin{figure}[htbp]
\centering
% FIGURE-SPEC-BEGIN
% ID: fig:packaging-artifact-flow
% Source of truth:
%   - ../.github/workflows/genesys-debian-package.yml
%   - ../scripts/run-debian-package-workflow.sh
%   - ../debian/control
%   - ../debian/genesys-common.install
%   - ../debian/genesys-gui.install
%   - ../debian/genesys-shell.install
%   - ../debian/genesys-worker.install
% Purpose:
%   Show the relationship between the Debian source package, the five binary
%   packages, and the generated diagnostics/package/lifecycle-evidence
%   artifact outputs across both workflow jobs.
% Required visual content:
%   - source package
%   - five binary packages
%   - genesys-debian-diagnostics artifact (Job 1)
%   - genesys-debian-packages artifact (Job 1, consumed by Job 2)
%   - genesys-debian-lifecycle-evidence artifact (Job 2)
% Accessibility description:
%   A package artifact map that separates source, build outputs, diagnostics,
%   and lifecycle evidence.
% Validation:
%   The package names and artifact names must match the current repository.
% FIGURE-SPEC-END
\figureplaceholder{figs/generated/packaging-artifact-flow}
\caption{Current artifact flow for Debian packaging.}
\label{fig:packaging-artifact-flow}
\end{figure}

Figure~\ref{fig:packaging-artifact-flow} should separate source, binary, and
diagnostic outputs so that the reader can see which artifacts are meant for
installation and which are only for inspection.

\begin{figure}[htbp]
\centering
% FIGURE-SPEC-BEGIN
% ID: fig:packaging-package-layout
% Source of truth:
%   - ../debian/control
%   - ../debian/rules
%   - ../debian/genesys-common.install
%   - ../debian/genesys-gui.install
%   - ../debian/genesys-shell.install
%   - ../debian/genesys-worker.install
%   - ../packaging/linux/io.github.rlcancian.genesys.desktop
%   - ../packaging/linux/io.github.rlcancian.genesys.metainfo.xml
% Purpose:
%   Show the installed filesystem layout across the five binary packages: the
%   public wrapper commands, the libexec launcher/dispatcher/system-fallback
%   tree, and the desktop/metainfo/icon files.
% Required visual content:
%   - /usr/bin/{genesys-gui,genesys-shell,genesys-worker,genesys-web,genesys-mcp}
%   - /usr/libexec/genesys/{genesys-launcher,genesys-dispatch}
%   - /usr/libexec/genesys/system/bin/{genesys-gui,genesys-shell,genesys-worker}
%   - /etc/genesys/update.conf
%   - desktop file, metainfo file, icon path
% Accessibility description:
%   A package contents diagram grouping public wrappers, the stable
%   launcher/dispatcher mechanism, the system fallback binaries, and desktop
%   integration assets by owning package.
% Validation:
%   The file paths and package names must match the current Debian manifests.
% FIGURE-SPEC-END
\figureplaceholder{figs/generated/packaging-package-layout}
\caption{Current Debian package contents and install layout.}
\label{fig:packaging-package-layout}
\end{figure}

Figure~\ref{fig:packaging-package-layout} should make it obvious that
\texttt{genesys-common} owns the stable launcher/dispatcher mechanism and
administrative policy, each application package owns its own public wrapper
and system fallback binary, and \texttt{genesys-web} owns no payload of its
own.

\section{The Launcher/Dispatcher packaging contract}
\label{sec:packaging-launcher-contract}

The public commands installed by these packages are deliberately tiny
wrappers, not the application binaries themselves. Each one execs a fixed,
absolute path with a fixed set of arguments, so its behavior is stable
independent of any future change to per-user runtime layout:
\begin{itemize}
\item \texttt{/usr/bin/genesys-gui} execs
  \texttt{/usr/libexec/genesys/genesys-launcher --app genesys-gui\allowbreak{} --interactive-update "\$@"};
\item \texttt{/usr/bin/genesys-shell}, \texttt{/usr/bin/genesys-worker},
  \texttt{/usr/bin/genesys-web}, and \texttt{/usr/bin/genesys-mcp} each exec
  \texttt{/usr/libexec/genesys/genesys-dispatch --app <name> -- "\$@"}.
\end{itemize}
These wrapper templates live at
\texttt{source/applications/launcher/wrappers/launcher-wrapper.in} and
\texttt{dispatch-wrapper.in} and are configured by CMake at build time, so the
generic CMake install tree and the Debian package install the identical
contract (Section~\ref{sec:packaging-cmake-install-tree}).

The launcher (\texttt{genesys-gui} only) and the dispatcher (every other
command) resolve which runtime actually serves the request through the same
\texttt{RuntimeSelector} contract validated by the Launcher's own focused
tests: a validated per-user runtime referenced by
\texttt{\textasciitilde/.local/share/genesys/current}, or the Debian-provided
system fallback under \texttt{/usr/libexec/genesys/system/bin/}. Only the
launcher may perform a network update check; the dispatcher never does.
\texttt{genesys-web} is resolved by the selector as a compatibility alias for
the \texttt{genesys-worker} application, not as a separate application.

\texttt{genesys-common} does not install a system fallback for
\texttt{genesys-mcp}: \texttt{source/applications/mcp/} is a Python
application whose dependencies (\texttt{mcp}, \texttt{httpx}) are outside the
Debian build contract. The Debian build does not fetch anything from PyPI,
vendor those dependencies, or reimplement MCP in C++. If no per-user runtime
provides \texttt{bin/genesys-mcp}, the dispatcher returns a controlled
diagnostic and exit status 127 instead of failing unpredictably.

\section{Generic CMake install tree and DESTDIR staging}
\label{sec:packaging-cmake-install-tree}

The Launcher, GUI, Shell, and Worker install rules use \texttt{GNUInstallDirs}
and are written once, in the generic CMake tree, independent of Debian. The
Debian package build (\texttt{debian/rules}) configures with
\texttt{-DCMAKE\_INSTALL\_PREFIX=/usr} and lets \texttt{dh\_auto\_install} run
\texttt{DESTDIR=debian/tmp ninja install} through debhelper, exactly like any
other \texttt{cmake+ninja} package. The \texttt{debian/*.install} files then
only partition that single staged tree between the five binary packages; they
never copy an executable directly from the build tree, and no Debian-specific
install logic exists outside \texttt{debian/rules} and the per-package
\texttt{.install}/\texttt{.manpages} files.

A maintainer can reproduce the same staged tree without \texttt{dpkg-buildpackage}
for a quick contract check:
\begin{lstlisting}[language=bash,caption={Staged install without touching the running system},label={lst:packaging-destdir-install}]
cmake --preset launcher-app
cmake --build --preset launcher-app --parallel "$(nproc)"
DESTDIR=/tmp/genesys-stage cmake --install build/launcher-app --prefix /usr
\end{lstlisting}
This is exactly the check the Launcher CI job performs after building; it
never writes into \texttt{/usr} on the runner itself.

\section{Package lifecycle validation}
\label{sec:packaging-lifecycle-validation}

\texttt{packaging/linux/validate-debian-lifecycle.sh} is the script executed
by the workflow's second job. It takes the built \texttt{.deb} directory and
an output directory as arguments, and exercises, in order, against the exact
five packages produced by Job 1:
\begin{enumerate}
\item package contents/inventory inspection and forbidden-artifact scanning
  (no \texttt{build-debian}, no \texttt{.git}, no stray \texttt{.deb}s inside
  a \texttt{.deb});
\item wrapper-contract checks (each public command execs the fixed
  launcher/dispatcher path with the fixed arguments, and contains no
  \texttt{\$HOME}/\texttt{\textasciitilde/} reference);
\item installation with \texttt{apt}, so dependency resolution runs
  normally, followed by \texttt{dpkg-query}/\texttt{dpkg -S} ownership checks
  for every public, libexec, and configuration path;
\item execution as a disposable, non-\texttt{sudoers} user
  (\texttt{genesysstudent}) with a clean home directory: system fallback for
  Shell, Worker (bounded, loopback-only \texttt{/health} on an ephemeral
  port), the missing-MCP controlled diagnostic, and a bounded Xvfb GUI start
  with \texttt{--no-update};
\item a fake, schema-valid per-user runtime (\texttt{VERSION},
  \texttt{MANIFEST.json}, \texttt{bin/*} fixtures) selected instead of the
  system fallback, with arguments preserved through the dispatcher;
\item the administrative \texttt{allow\_user\_runtime=false} override forcing
  system fallback even with a valid \texttt{current} runtime present;
\item six invalid-runtime cases (missing/broken \texttt{current}, missing
  \texttt{VERSION}, missing \texttt{MANIFEST.json}, missing or
  non-executable application binary), each expected to fall back cleanly with
  no crash and no network access;
\item package-owned file integrity (checksums and \texttt{dpkg -V}) before and
  after all per-user-runtime exercises;
\item reinstall with a locally edited conffile, confirming the local edit
  survives and the active per-user runtime is untouched;
\item \texttt{remove}, confirming package-owned paths are gone, the conffile
  and per-user runtime remain;
\item \texttt{purge}, confirming the conffile is gone and the per-user
  runtime still remains.
\end{enumerate}
No step downloads a GenESyS runtime from GitHub Releases or otherwise reaches
outside the packages produced by the same workflow run. Section
\ref{sec:packaging-release-boundary} covers what this validation does and does
not authorize.

\section{Release boundary and limitations}
\label{sec:packaging-release-boundary}

The current automation surface still has clear boundaries:
\begin{itemize}
\item the ordinary CI workflow and the Debian package workflow are distinct;
\item the Debian workflow now triggers on \texttt{WorkInProgress},
  \texttt{20262}, \texttt{currentStable}, and \texttt{master}, matching the
  ordinary CI trigger surface for packaging-relevant path changes -- it is no
  longer filtered to a dated \texttt{WiP20261}/\texttt{20261} pair;
\item within one workflow run, a green build job does unlock install,
  non-root-user, system-fallback, per-user-runtime, reinstall, remove, and
  purge validation through the dependent lifecycle job
  (Section~\ref{sec:packaging-lifecycle-validation}) -- but a green build job
  alone, without a green lifecycle job on the same run, still does not
  constitute that evidence;
\item AppStream and lintian checks validate package metadata quality, not the
  supported-feature boundary of the whole project;
\item a \texttt{genesys-debian-packages} GitHub Actions artifact is a
  validation artifact for inspection and manual testing. It is not a GitHub
  Release, it is not published to any APT/PPA repository, and package build
  or lifecycle success does not by itself authorize public release,
  official package signing, or PPA publication;
\item no runtime signing key, GitHub Release, or APT/PPA publication is
  produced by this workflow; those remain separate, explicitly authorized
  release activities.
\end{itemize}

The chapter therefore treats packaging as a validated build, install, and
lifecycle surface, not as a release decision.

\begin{figure}[htbp]
\centering
% FIGURE-SPEC-BEGIN
% ID: fig:packaging-release-boundary
% Source of truth:
%   - ../.github/workflows/genesys-ci.yml
%   - ../.github/workflows/genesys-debian-package.yml
%   - ../docs/ai_assistants/GOVERNANCE.md
%   - ../docs/ai_assistants/STATUS.md
%   - ../docs/ai_assistants/BACKLOG_HUMAN.md
% Purpose:
%   Show the boundary between CI/package validation and release authorization,
%   including the fact that package success is not a public release decision.
% Required visual content:
%   - CI validation box
%   - Debian package validation box
%   - metadata checks
%   - release or publication gate
%   - human approval boundary
% Accessibility description:
%   A gate diagram that distinguishes technical validation from release approval.
% Validation:
%   The diagram must not suggest that package build success alone authorizes
%   release.
% FIGURE-SPEC-END
\figureplaceholder{figs/generated/packaging-release-boundary}
\caption{Current release boundary for packaging and CI.}
\label{fig:packaging-release-boundary}
\end{figure}

Figure~\ref{fig:packaging-release-boundary} should make the release boundary
explicit. Technical validation is necessary, but it is not the same thing as
publication authority.
